\textbf{Batch reservation cost.}
All $M$ VMs are acquired at time zero and remain reserved until the batch completes at $H_{\mathrm{max}}(s)$. A VM that finishes its own queue early, or receives no task, is still charged until this common release. There is no separate VM release, provisioning delay, minimum charge, or billing-time rounding. Let $r_j$ be the synthetic rate in cost units per second. The batch reservation cost is
\begin{equation}\label{eq:cost}
\begin{aligned}
C(s)&=H_{\mathrm{max}}(s)\sum_{j=1}^{M}r_j\\
&=\sum_{j=1}^{M}r_jH_j(s)+\sum_{j=1}^{M}r_j[H_{\mathrm{max}}(s)-H_j(s)].
\end{aligned}
\end{equation}
The two terms on the second line charge busy time and reserved idle time. This is a fixed-fleet, continuous-time adaptation of reservation charging~\cite{Oprescu2010}, rather than the full BaTS model or ERTH's resource-bandwidth cost~\cite{Qin2024}.

For configured positive processing-rate bounds $\nu_{\mathrm{min}}\leq\nu_j\leq\nu_{\mathrm{max}}$, a denominator floor $\epsilon_{\nu}>0$ in MIPS, and $0\leq r_{\mathrm{min}}\leq r_{\mathrm{max}}$, the retained price mapping is
\begin{equation}\label{eq:rate}
r_j=r_{\mathrm{min}}+(r_{\mathrm{max}}-r_{\mathrm{min}})\sqrt{\frac{\nu_j-\nu_{\mathrm{min}}}{\max\{\nu_{\mathrm{max}}-\nu_{\mathrm{min}},\epsilon_{\nu}\}}}.
\end{equation}
The bounds belong to the configured capacity profile. The mapping is nondecreasing but not proportional to capacity, and it is synthetic rather than a calibrated cloud tariff. Cost depends on the batch horizon and fleet rates, not on per-task FCFS waiting. The earlier completion-time proxy $\sum_i r_{j(i)}CT_i$ is retained only as an exploratory control and is not used in the revised tables or figures.
